\documentclass[11pt]{article}
\usepackage[margin=2.5cm]{geometry}
\usepackage{amsmath,amssymb}
\usepackage{algorithm}
\usepackage{algpseudocode}
\usepackage{hyperref}

\title{PATHX -- Implementation of GALOIS Features}
\author{Minseok Kim \and Juliette Tonneau \and Danil Smirnov \and Bita Khashechian}
\date{5 December 2025}

\begin{document}
\maketitle

\tableofcontents

\section{Introduction}

The GALOIS framework proposes to treat large language models (LLMs) as
\emph{logical operators} inside a relational query engine, instead of asking them to
directly answer a query in a single natural language pass. Rather than hoping that one
prompt produces the full and correct result, GALOIS decomposes a query into smaller,
structured steps (such as scanning keys, filling tuples, or enforcing predicates) and uses
the LLM as a component inside an execution plan.

In this report we focus on reproducing the key ideas of GALOIS in our own system
and on applying them to the same EU Cities example as in the paper. We start from
simple baselines, where the LLM is asked to respond in one shot, and then move to
more sophisticated variants that mimic database-style plans:

\begin{itemize}
  \item a basic non-iterative baseline (single-pass NL or SQL prompting),
  \item iterative baselines that re-ask the LLM for more rows,
  \item GALOIS variants wO, A, S, and F that structure the interaction via operators.
\end{itemize}

The goal is to understand how much the query \emph{structure} and the choice of plan
influence the quality of the tuples retrieved by the model. Throughout the report we
compare our implementations and numerical results to those described in the original
GALOIS paper, with particular attention to the EU Cities query.

\section{Baselines Implementation}
\label{sec:baselines}

\subsection{Non-iterative baselines}

The non-iterative baselines correspond to the most direct way of using an LLM: a single
prompt, a single answer, and no further refinement. We implemented two such baselines.

In the natural-language (NL) baseline, the model receives a description of the
dataset and the task in plain language, together with an instruction to return the answer
as a table in JSON format. The system sends the prompt once and parses the resulting
JSON into a set of tuples. No further calls are made, and no consistency checks are
performed.

In the SQL baseline, the model is asked to produce a SQL query that answers the
task. The generated SQL is then executed on the ground-truth database. This baseline
often performs better than the pure NL version because the model is guided towards
structured queries, but it still relies on the ability of the LLM to write correct SQL in a
single step.

These baselines are simple and cheap in terms of token usage. However, they are
also brittle: when the query involves several constraints or multiple attributes, the model
tends to ignore part of the conditions or to produce incomplete result sets. Since the
system does not revisit the answer, those errors propagate unchanged to the final output.

\subsection{Iterative baselines}

To address the limitations of single-pass prompting, we implemented iterative versions of
the baselines. The motivation is empirical: when we ask an LLM to list entities (cities,
movies, airlines, \dots), the first response is rarely complete. The model often lists a few
prominent rows and stops, even if the prompt explicitly says ``list all''.

The iterative baselines therefore repeat the interaction. The first prompt describes
the dataset, the table, the attributes, and the filtering conditions, and asks the model
to return a JSON table. The system parses the output, stores the rows, and constructs
a second prompt that includes the previously retrieved rows together with the instruction
not to repeat them. This process is repeated for a fixed number of iterations or until the
model stops producing new rows.

In the iterative NL baseline, all of this happens in natural language. In the iterative
SQL baseline, the model is instead asked to refine or extend its previous SQL statement.
In both cases, we observed that iteration improves recall substantially, but the model
still struggles to systematically enforce all constraints. This motivates moving from
``iterated prompting'' to a more principled, operator-based design inspired by GALOIS.

\section{Architecture and Implementation Details}
\label{sec:architecture}

\subsection{Overview}

Our implementation follows a layered architecture that separates:

\begin{itemize}
  \item \emph{query handling}, where a SQL query is received and parsed;
  \item \emph{planning}, where we decide which GALOIS variant and which operators to use;
  \item \emph{optimization}, where we choose which WHERE conditions to push down;
  \item \emph{iterative execution}, where LLMScan operators interact with the LLM.
\end{itemize}

At the top, a query interface takes a SQL query and maps it to one of several execution
plans: a baseline, GALOIS wO, A, S, or F. The core of the system is an implementation
of the LLM-based scan operators described in the GALOIS paper, together with a
pushdown optimizer that decides which predicates should be enforced inside the LLM
and which should be enforced afterwards on the retrieved tuples.

The architecture is designed so that all variants share the same building blocks:

\begin{itemize}
  \item a prompt layer, which defines the templates used to talk to the LLM,
  \item scan operators (Table-Scan and Key-Scan) that implement LLMScan,
  \item a small logical optimizer for predicates,
  \item and a thin experimental layer that prepares ground truth and evaluates metrics.
\end{itemize}

In the following subsections we describe each of these components and how they relate
to the algorithms in the GALOIS paper.

\subsection{Prompt layer and LLM interaction}

All interactions with the model are mediated through a small prompt layer built with
Jinja2 templates. Instead of constructing prompts ad-hoc, we define one template per
type of operator or evaluation step. The most important templates are:

\begin{itemize}
  \item \emph{Table-Scan prompts}, which ask the model to list rows from a given table,
    restricted to specific attributes and an optional WHERE clause.
  \item \emph{Key-Scan prompts}, which ask only for key attributes (for example the city
    name or a primary key column).
  \item \emph{Tuple-completion prompts}, which, given a particular key value, ask the
    model to return all remaining attributes for that row.
  \item \emph{Cell-level prompts}, which ask for a single attribute value for one key.
    These are used only in the wO (without optimizations) variant.
  \item \emph{Condition-evaluation prompts}, which present each atomic WHERE condition
    and ask the model to label it as either ``higher''/``lower'' (selectivity mode) or
    ``high''/``low'' (confidence mode).
  \item \emph{Key-confidence prompts}, which ask for a scalar in $[0,1]$ describing how
    confident the model is that it can enumerate keys satisfying the query.
\end{itemize}

All templates enforce JSON output. A shared base class, \texttt{BaseScanOperator},
wraps the LLM client and provides a helper to send a prompt, measure wall-clock time,
read token usage, and parse the JSON. This ensures that Table-Scan, Key-Scan, and the
wO variant all use the same interface to the LLM and that timing and token statistics
are collected consistently.

\subsection{LLMScan as physical operators}

In the terminology of the GALOIS paper, LLMScan is the logical operator that retrieves
tuples from the model. In our implementation, LLMScan is realized by two physical
operators that inherit from \texttt{BaseScanOperator}:

\begin{itemize}
  \item \textbf{TableScanOperator}, which retrieves full tuples directly,
  \item \textbf{KeyScanOperator}, which first retrieves keys and then completes tuples.
\end{itemize}

Both operators are iterative. The first iteration asks the model for an initial batch of
tuples or key values. Subsequent iterations include the previously returned rows in the
prompt and explicitly instruct the model not to repeat them. After each call, rows are
deduplicated via a row-signature function before being added to the global result set.

\subsubsection{Table-Scan}

Table-Scan corresponds to Algorithm~1 in the GALOIS paper. The operator takes as
input a dataset name, a table name, a list of projected attributes, an optional WHERE
clause, and a maximum number of iterations. It then repeatedly calls the LLM to list
rows that satisfy the query.

\begin{algorithm}[h]
  \caption{Table-Scan}
  \label{alg:tablescan}
  \begin{algorithmic}[1]
    \Require dataset, table, attributes, where\_clause, max\_iters
    \State $R \gets \emptyset$ \Comment{collected rows}
    \State $p \gets \text{build\_first\_table\_prompt}(\text{dataset}, \text{table}, \text{attributes}, \text{where\_clause})$
    \State $rows \gets \text{call\_LLM}(p)$
    \State $R \gets R \cup rows$
    \For{$i \gets 2$ to $max\_iters$}
      \State $p \gets \text{build\_followup\_table\_prompt}(\text{dataset}, \text{table}, \text{attributes}, \text{where\_clause}, R)$
      \State $rows \gets \text{call\_LLM}(p)$
      \State $rows \gets rows \setminus R$ \Comment{keep only new rows}
      \If{$rows = \emptyset$}
        \State \textbf{break}
      \EndIf
      \State $R \gets R \cup rows$
    \EndFor
    \State \Return $R$
  \end{algorithmic}
\end{algorithm}

This operator is token-efficient when the model can produce many rows per call. For
simple queries, the first one or two iterations already cover most of the relevant tuples.
For more complex queries, however, it may struggle to satisfy all constraints inside a
single generative step, which motivates introducing Key-Scan.

\subsubsection{Key-Scan}

Key-Scan corresponds to Algorithm~2 in the GALOIS paper. It decomposes retrieval into
two phases: key discovery and tuple completion. In the first phase, the operator asks the
LLM to list only key attributes. In the second phase, it iterates over these keys and uses
more focused prompts to fill in the remaining attributes for each key.

\begin{algorithm}[h]
  \caption{Key-Scan}
  \label{alg:keyscan}
  \begin{algorithmic}[1]
    \Require dataset, table, attributes, key\_attributes, where\_clause, max\_iters
    \State Split \textit{attributes} into key attributes $K$ and non-key attributes $A$
    \State $Keys \gets \emptyset$
    \Comment{Phase 1: key discovery}
    \State $p \gets \text{build\_first\_key\_prompt}(\text{dataset}, \text{table}, K, \text{where\_clause})$
    \State $new \gets \text{call\_LLM}(p)$
    \State $Keys \gets Keys \cup new$
    \For{$i \gets 2$ to $max\_iters$}
      \State $p \gets \text{build\_followup\_key\_prompt}(\text{dataset}, \text{table}, K, \text{where\_clause}, Keys)$
      \State $new \gets \text{call\_LLM}(p)$
      \State $new \gets new \setminus Keys$
      \If{$new = \emptyset$}
        \State \textbf{break}
      \EndIf
      \State $Keys \gets Keys \cup new$
    \EndFor
    \Comment{Phase 2: tuple completion}
    \State $R \gets \emptyset$
    \ForAll{$k \in Keys$}
      \State $p \gets \text{build\_tuple\_prompt}(\text{dataset}, \text{table}, k, A, \text{where\_clause})$
      \State $vals \gets \text{call\_LLM}(p)$
      \State $row \gets k \cup vals$
      \State $R \gets R \cup \{row\}$
    \EndFor
    \State \Return $R$
  \end{algorithmic}
\end{algorithm}

This two-phase strategy typically yields higher-quality results because the prompts in
the second phase are more focused. The trade-off is that we need at least one LLM call
per key to complete the tuple, which can increase cost when many keys are present.

\subsection{Predicate pushdown and logical optimization}

To decide which parts of the WHERE clause should be enforced inside the LLM and
which should be enforced afterwards, we implement a small logical optimizer.

First, the WHERE clause is parsed into a list of atomic conditions separated by
\texttt{AND}. A separate component calls the LLM with a condition-evaluation prompt
and obtains a map from conditions to labels:

\begin{itemize}
  \item in \emph{selectivity mode}, each condition is labeled ``higher'' or ``lower'';
  \item in \emph{confidence mode}, each condition is labeled ``high'' or ``low''.
\end{itemize}

Given these labels, the optimizer applies the rule described in the GALOIS paper. Let
$C$ be the set of conditions and $C^\ast$ the subset whose label matches the chosen
criterion (``higher'' or ``high''):

\begin{itemize}
  \item if $|C^\ast| = 0$, no condition is pushed down;
  \item if $|C^\ast| = 1$, that single condition is pushed down;
  \item if $|C^\ast| > 1$, all conditions in $C^\ast$ are pushed down and combined with
        \texttt{AND}.
\end{itemize}

Conditions that are not pushed down are treated as residual predicates. After
LLMScan finishes, the system re-applies the full original WHERE clause on the retrieved
rows using DuckDB. In this way, even if the LLM misinterprets the prompt or partially
ignores some conditions, the final filtering step recovers the correct semantics.

\subsection{GALOIS variants wO, A, S, F}

On top of the scan operators and the pushdown optimizer, we implement the four
GALOIS variants described in the paper. Each variant takes a dataset, table, attribute
list, and WHERE clause, and produces a result set and metadata.

All variants share a common skeleton: they extend the attribute list to include any
columns mentioned in the WHERE clause, run one or more scan operators, re-apply the
full WHERE clause locally, and finally project the tuples back to the attributes requested
by the original query.

\paragraph{Galois wO (without optimizations).}

This variant behaves like a cell-level LLM oracle. It first discovers keys using a key-only
scan, then for each key and each non-key attribute it calls the cell-level prompt to fill in
one value. No WHERE pushdown is performed; all filtering happens after the tuples are
assembled. The number of LLM calls is therefore proportional to the number of rows
times the number of attributes. This variant is deliberately expensive and serves as a
reference for what is possible without any logical or physical optimizations.

\paragraph{Galois A.}

Galois A uses Table-Scan and pushes the entire WHERE clause into the prompt. After
LLMScan, the same WHERE clause is re-applied locally. This variant is cheap in terms
of LLM calls but can miss valid tuples if the model struggles to satisfy every condition
simultaneously. It corresponds to a simple ``push everything'' strategy.

\paragraph{Galois S.}

Galois S uses the selectivity labels from the LLM. Conditions labeled as ``higher'' are
candidates for pushdown; conditions labeled ``lower'' are left as residual predicates.
The optimizer applies the 0/1/$>$1 rule: if exactly one condition is ``higher'', it is
pushed; if several are ``higher'', they are all pushed; if none are ``higher'', nothing is
pushed. Table-Scan is used as the physical operator. In practice, this variant often
preserves recall compared to Galois A because it avoids pushing down low-selectivity
conditions that might cut off relevant tuples too early.

\paragraph{Galois F.}

Galois F is the full optimizer. For the logical part, it uses confidence labels
(``high''/``low'') in exactly the same way as Galois S uses selectivity labels: conditions
marked ``high'' are considered for pushdown, and the 0/1/$>$1 policy is applied.

For the physical part, it asks the LLM for a key-retrieval confidence $c \in [0,1]$.
If the query projects $n$ attributes, we compute
\[
  \mathrm{conf}(q) = c^{n},
\]
following the idea in the paper of normalizing confidence by query width. This
$\mathrm{conf}(q)$ is compared to a threshold $\tau$ (tunable, default $0.6$). If
$\mathrm{conf}(q) \ge \tau$, the system chooses Key-Scan; otherwise it chooses
Table-Scan. In this way, Galois F combines predicate pushdown with a dynamic choice
between the two physical operators.

In all cases, the final stage is to re-apply the full WHERE clause locally and project the
rows back to the attributes requested by the original query.

\section{Experimental Setup and EU Cities Query}
\label{sec:eu_cities_setup}

\subsection{Data and ground truth}

To evaluate our implementation in a way that is comparable to the GALOIS paper, we
construct ground truth by executing the official SQL queries on the benchmark datasets
in DuckDB. A dedicated script, \texttt{build\_ground\_truth.py}, reads the CSV files for
each dataset, creates tables in a DuckDB database, and runs each \texttt{SELECT} query
from the corresponding \texttt{queries\_*.sql} file. The result of each query is stored as
\texttt{query<i>.json} under a dataset-specific directory.

These JSON files serve as the gold-standard answers for all methods: baselines and
GALOIS variants are evaluated against them using the same metrics (F1-Cell, Cardinality,
Tuple-Constraint, AVG).

\subsection{Pipeline for baselines and GALOIS}

The experimental pipeline follows the same broad structure for every run:

\begin{enumerate}
  \item \textbf{Ground truth construction.} Run \texttt{build\_ground\_truth.py} on the
    target datasets to produce the reference JSON files.
  \item \textbf{Dataset ingestion.} Load the same CSV files into DuckDB so that SQL
    queries can be executed directly when needed.
  \item \textbf{Baseline execution.} Run the NL and SQL baselines (non-iterative and
    iterative variants) and save their predicted result sets in a submissions directory.
  \item \textbf{GALOIS execution.} Run the GALOIS query engine with each of the
    variants (wO, A, S, F), using a fixed LLM provider and configuration, and save
    the outputs and metadata (tokens, time).
  \item \textbf{Evaluation.} Compare each submission against the ground truth using
    the GALOIS metrics and aggregate results by dataset and overall.
\end{enumerate}

Evaluation is handled by the script \texttt{galois\_eval.py}, which reads the ground-truth
and submission JSON files, aligns them by query, and computes F1-Cell, Cardinality,
Tuple-Constraint, and AVG. When metadata is present, it also reports token usage and
average response time.

\subsection{EU Cities query}

To make our results directly comparable to the GALOIS paper, we use the same EU
Cities example. The \texttt{EU\_Cities} table is part of the WORLD dataset and includes
European cities with attributes such as \texttt{name}, \texttt{country}, \texttt{population},
and \texttt{num\_private\_hospitals}. The target query selects all cities with more than
1\,000\,000 inhabitants and more than 15 private hospitals, and can be written as:

\begin{verbatim}
SELECT name, country, population, num_private_hospitals
FROM EU_Cities
WHERE population > 1000000
  AND num_private_hospitals > 15;
\end{verbatim}

We run this query in five ways:

\begin{itemize}
  \item a naive single-shot NL baseline,
  \item an iterative NL baseline,
  \item Galois wO (cell-by-cell, no pushdown),
  \item Galois A (Table-Scan with full pushdown),
  \item Galois S (Table-Scan with selectivity-based pushdown),
  \item Galois F (full optimization with confidence-based pushdown and dynamic choice
        between Table-Scan and Key-Scan).
\end{itemize}

For each execution plan, we record the tuples returned, the number of tokens consumed,
and the total execution time. These are then compared against the exact SQL result for
the EU Cities query using the GALOIS metrics. The detailed numerical results and the
corresponding tables of returned cities are reported in the Results section of the report.

\section{Results}

[Here you keep your existing Tables for NL, SQL, GALOIS~wO, A, S, F, including the
EU Cities tables and aggregate metrics across datasets. You can add one short bridging
paragraph such as: “We now summarize the numerical results obtained by the pipelines
described in Sections~\ref{sec:baselines}, \ref{sec:architecture}, and
\ref{sec:eu_cities_setup}.”]

\section{Discussion}

[Keep your original discussion text. It already matches the story told by the
implementation: baselines vs.\ operator-based plans, behavior of each variant, and the
trade-offs between cost and completeness.]

\section{Conclusion}

[Keep your original conclusion text. No changes needed here; it remains consistent with
the architecture and experiments described above.]

\end{document}
